\section*{अध्याय \ref{ch:g10:cell-metabolism} --- कोशिका का उपापचय}

\begin{solution}{exo:g10:cell-metabolism:1}
\omterm{def:g10:cell-metabolism:metabolism}{उपापचय} किसी \omterm{def:g10:cells-common-unit:cell}{कोशिका} की सारी रासायनिक अभिक्रियाओं का समुच्चय है, जो ऊर्जा
के लिए अणु तोड़ती हैं और \omterm{def:g10:cells-common-unit:cell}{कोशिका} के अपने अणु बनाती हैं। \omterm{def:g10:cell-metabolism:metabolism}{एंजाइम} वह \omterm{def:g10:chemistry-of-life:families}{प्रोटीन}
है जो किसी एक विशिष्ट अभिक्रिया को तेज़ करता है, अपने क्रियाधार से जुड़कर
उत्पाद छोड़ता है, और ख़ुद ख़र्च नहीं होता।
\end{solution}

\begin{solution}{exo:g10:cell-metabolism:2}
श्वसन: $\mathrm{C_6H_{12}O_6} + 6\,\mathrm{O_2} \to
6\,\mathrm{CO_2} + 6\,\mathrm{H_2O}$ + ऊर्जा। \omterm{def:g10:cell-metabolism:autotrophy}{प्रकाशसंश्लेषण}:
$6\,\mathrm{CO_2} + 6\,\mathrm{H_2O}$ + प्रकाश $\to \mathrm{C_6H_{12}O_6}
+ 6\,\mathrm{O_2}$। पदार्थ का कुल हिसाब उल्टा है; पर अभिक्रियाएँ,
\omterm{def:g10:cell-metabolism:metabolism}{एंजाइम}, जगहें (\omterm{def:g10:cells-common-unit:organelle}{सूत्रकणिका} बनाम \omterm{def:g10:cells-common-unit:organelle}{हरितलवक}) और ऊर्जा-स्रोत (रासायनिक बंध
बनाम प्रकाश) पूरी तरह अलग हैं।
\end{solution}

\begin{solution}{exo:g10:cell-metabolism:3}
\omterm{def:g10:cell-metabolism:autotrophy}{परपोषी}: कुकुरमुत्ता, गाजर की जड़ की \omterm{def:g10:cells-common-unit:cell}{कोशिका}, दही का जीवाणु, मनुष्य की
पेशी \omterm{def:g10:cells-common-unit:cell}{कोशिका}। \omterm{def:g10:cell-metabolism:autotrophy}{स्वपोषी}: मॉस, क्लोरेला।
\end{solution}

\begin{solution}{exo:g10:cell-metabolism:4}
गुब्बारे को कार्बन डाइऑक्साइड फुलाती है; और तरल में इथेनॉल जमा होता है।
फ़्लास्क में ऑक्सीजन की कमी है: हवा के बुलबुले छोड़े जाते तो खमीर श्वसन
करता, इथेनॉल नहीं बनाता और चीनी कहीं कम ख़र्च करता।
\end{solution}

\begin{solution}{exo:g10:cell-metabolism:5}
कैटेलेज़ एक \omterm{def:g10:chemistry-of-life:families}{प्रोटीन} है; उबालने से उसका आकार और इसलिए उसकी सक्रियता नष्ट
हो जाती है, और हाइड्रोजन परॉक्साइड अपने आप तेज़ी से नहीं टूटता।
\end{solution}

\begin{solution}{exo:g10:cell-metabolism:6}
$(7.9 - 0.9)/7 = \qty{1.0}{mg/L}$ प्रति मिनट। वक्र इसलिए चपटा पड़ता है
कि ऑक्सीजन चुक जाती है; खमीर तब \omterm{def:g10:cell-metabolism:fermentation}{किण्वन} पर चला जाता है, जो कोई ऑक्सीजन
नहीं खपाता।
\end{solution}

\begin{solution}{exo:g10:cell-metabolism:7}
$0.8 + 0.2 = \qty{1.0}{mg/L}$ प्रति मिनट: प्रकाश वाली ढाल उत्पादन में से
प्रकाश में भी चलते रहने वाले श्वसन को घटाकर बची ढाल है।
\end{solution}

\begin{solution}{exo:g10:cell-metabolism:8}
\qty{18}{g} यानी \qty{0.1}{mol} ग्लूकोज़; उसे \qty{0.6}{mol} ऑक्सीजन
चाहिए, यानी $0.6 \times 32 = \qty{19.2}{g}$, और वह \qty{0.6}{mol} कार्बन
डाइऑक्साइड छोड़ता है, यानी $0.6 \times 44 = \qty{26.4}{g}$।
\end{solution}

\begin{solution}{exo:g10:cell-metabolism:9}
अंकुरित होते मटर अपने भंडारों का श्वसन करते हैं और ऊष्मा छोड़ते हैं; उबले
मटर मरे हुए हैं और उनके \omterm{def:g10:cell-metabolism:metabolism}{एंजाइम} नष्ट हैं, इसलिए न अभिक्रिया, न ऊष्मा।
हवाबंद कर देने पर ऑक्सीजन चुक जाती और श्वसन रुक जाता।
\end{solution}

\begin{solution}{exo:g10:cell-metabolism:10}
खुला आधा हिस्सा नीला-काला रंग पकड़ता है (प्रकाश में \omterm{def:g10:cell-metabolism:autotrophy}{प्रकाशसंश्लेषण} से बना
मंड); ढँका हुआ आधा नहीं। 48 घंटे का अँधेरा पत्ते को उसका पहले वाला मंड
ख़र्च कर लेने देता है, ताकि जो मंड मिले वह उसी दिन के प्रकाश का हो।
\end{solution}

\begin{solution}{exo:g10:cell-metabolism:11}
\omterm{def:g10:cell-metabolism:fermentation}{किण्वन} प्रति ग्लूकोज़ श्वसन से लगभग पंद्रह गुना कम ATP देता है; इसलिए
अपनी ज़रूरत भर ATP पाने के लिए \omterm{def:g10:cell-metabolism:fermentation}{किण्वन} करती \omterm{def:g10:cells-common-unit:cell}{कोशिका} को प्रति मिनट पंद्रह
गुना अधिक चीनी तोड़नी पड़ती है।
\end{solution}

\begin{solution}{exo:g10:cell-metabolism:12}
दौड़ के दौरान पेशी की ATP की माँग उस ऑक्सीजन से आगे निकल गई जो रक्त
पहुँचा सकता था; \omterm{def:g10:cells-common-unit:cell}{कोशिकाएँ} लैक्टिक \omterm{def:g10:cell-metabolism:fermentation}{किण्वन} पर चली गईं, और उसका उत्पाद, यानी
लैक्टिक अम्ल, जमा होता गया। धीमी जॉगिंग में ऑक्सीजन की आपूर्ति साथ बनी
रहती है, अकेला श्वसन ही ATP दे देता है, और कोई अम्ल नहीं बनता।
\end{solution}

\begin{solution}{exo:g10:cell-metabolism:13}
पर्यावरण बदला (कोई प्रकाश नहीं, कार्बनिक भोजन); \omterm{def:g10:universal-dna:gene}{जीन} नहीं बदले। वंशज
\omterm{def:g10:cells-common-unit:organelle}{हरितलवक} दोबारा बना लेते हैं, इसलिए निर्देश कभी खोए ही नहीं थे: जो खोया वह
इस्तेमाल में आ रहा साज़-सामान था, न कि \omterm{def:g10:universal-dna:gene}{जीनों} की दी हुई संभावनाएँ।
\end{solution}

\begin{solution}{exo:g10:cell-metabolism:14}
प्रकाश में पौधा \omterm{def:g10:cell-metabolism:autotrophy}{प्रकाशसंश्लेषण} करता है और ऑक्सीजन तथा \omterm{def:g10:chemistry-of-life:mineral-organic}{कार्बनिक पदार्थ}
देता है (पत्तियाँ, जो घोंघा खाता है); घोंघा श्वसन करता है और पौधे को
कार्बन डाइऑक्साइड देता है। अँधेरे में पौधा केवल श्वसन करता है: दोनों
ऑक्सीजन खपाते हैं, कोई उसे बनाता नहीं, और दोनों कुछ ही दिनों में मर जाते
हैं।
\end{solution}

\begin{solution}{exo:g10:cell-metabolism:15}
प्रति मिनट $30/2 = 15$ गुना अधिक ग्लूकोज़। शराब बनाने वालों को इथेनॉल
चाहिए, जो केवल \omterm{def:g10:cell-metabolism:fermentation}{किण्वन} बनाता है; श्वसन करता खमीर शर्करा को कार्बन
डाइऑक्साइड, जल और और खमीर में बदल देता, ऐल्कोहॉल में नहीं।
\end{solution}

\begin{solution}{pb:g10:cell-metabolism:1}
\textbf{1.} ऐल्कोहॉली \omterm{def:g10:cell-metabolism:fermentation}{किण्वन}, क्योंकि बंद टंकी में ऑक्सीजन नहीं है:
$\mathrm{C_6H_{12}O_6} \to 2\,\mathrm{C_2H_5OH} +
2\,\mathrm{CO_2}$।

\textbf{2.} $200/180 \approx \qty{1.11}{mol}$।

\textbf{3.} इथेनॉल $2 \times 1.11 = \qty{2.22}{mol}$, यानी $2.22
\times 46 \approx \qty{102}{g}$; आयतन $102/0.79 \approx
\qty{129}{mL}$; यानी आयतन के हिसाब से लगभग 13\%।

\textbf{4.} कार्बन डाइऑक्साइड $2.22 \times 44 \approx \qty{98}{g}$;
गैस के रूप में $2.22 \times 24 \approx \qty{53}{L}$।

\textbf{5.} लगभग \qty{53000}{L} $= \qty{53}{m^3}$। तहख़ाना
$4 \times 6 \times 2.5 = \qty{60}{m^3}$ का है: छूटी हुई गैस उसे लगभग भर
ही देती।

\textbf{6.} $\mathrm{C_6H_{12}O_6} + 6\,\mathrm{O_2} \to
6\,\mathrm{CO_2} + 6\,\mathrm{H_2O}$।

\textbf{7.} ऑक्सीजन $6 \times 1.11 = \qty{6.67}{mol}$, यानी $6.67 \times 32
\approx \qty{213}{g}$; कार्बन डाइऑक्साइड $6.67 \times 44 \approx
\qty{293}{g}$।

\textbf{8.} तीन गुना अधिक: श्वसन ग्लूकोज़ के छहों कार्बन कार्बन
डाइऑक्साइड में बदल देता है, \omterm{def:g10:cell-metabolism:fermentation}{किण्वन} केवल दो, और बाक़ी चार इथेनॉल में बने
रहते हैं।

\textbf{9.} श्वसन के साथ पंद्रह गुना अधिक देर: हर ग्लूकोज़ पंद्रह गुना
अधिक ATP देता है।

\textbf{10.} इथेनॉल केवल \omterm{def:g10:cell-metabolism:fermentation}{किण्वन} बनाता है; हवा भरी टंकी कार्बन
डाइऑक्साइड, जल और खमीर देती, शराब नहीं।

\textbf{11.} प्रति लीटर $1.11 \times 70 \approx \qty{78}{kJ}$; और पूरी
टंकी के लिए \qty{78000}{kJ}, यानी \qty{78}{MJ}।

\textbf{12.} $78/4.2 \approx \qty{18.5}{\celsius}$।

\textbf{13.} $20 + 18.5 \approx \qty{38.5}{\celsius}$: यानी ठीक खमीर की
सीमा के किनारे, और गरम तहख़ाना उसे पार करा देता। बड़ी टंकी ऊष्मा धीरे
खोती है, क्योंकि उसके आयतन के मुक़ाबले उसकी सतह छोटी है, इसलिए उसे
जानबूझकर ठंडा करना पड़ता है।

\textbf{14.} \omterm{def:g10:cell-metabolism:metabolism}{एंजाइम} \omterm{def:g10:chemistry-of-life:families}{प्रोटीन} हैं; ऊष्मा उन्हें खोल देती है और वे काम करना
बंद कर देते हैं। अपने एंजाइमों के बिना खमीर शर्करा को बदल नहीं सकता,
चाहे वह कितनी ही बची हो।

\textbf{15.} इथेनॉल के रासायनिक बंधों में: प्रति मोल ग्लूकोज़ लगभग
\qty{2800}{kJ}, और इसीलिए इथेनॉल जलता है।

\textbf{16.} फ़र्श पर, एक परत में, जो \omterm{def:g10:cell-metabolism:fermentation}{किण्वन} के साथ मोटी होती जाती है।
नीचे तक ले जाई गई मोमबत्ती वहाँ बुझ जाती है जहाँ कार्बन डाइऑक्साइड ने
ऑक्सीजन को हटा दिया है --- यानी किसी आदमी के चपेट में आने से पहले की
चेतावनी।

\textbf{17.} लैक्टिक \omterm{def:g10:cell-metabolism:fermentation}{किण्वन} पर। वह पंद्रह गुना कम ATP देता है; मस्तिष्क
और हृदय अपनी ज़रूरतें उस तरह पूरी नहीं कर सकते, और लैक्टिक अम्ल जमा होता
जाता है। एक-दो मिनट में ही होश जाता रहता है।

\textbf{18.} अधिक सांद्रता पर इथेनॉल ख़ुद खमीर की झिल्लियों और एंजाइमों
को नुक़सान पहुँचाता है: \omterm{def:g10:cells-common-unit:cell}{कोशिकाओं} का अपना ही उत्पाद उनके \omterm{def:g10:cell-metabolism:metabolism}{उपापचय} को विष दे
देता है।

\textbf{19.} बुलबुलों में \omterm{def:g10:cell-metabolism:fermentation}{किण्वन} की कार्बन डाइऑक्साइड है। इथेनॉल भट्ठी
में उड़ जाता है (वह \qty{78}{\celsius} पर उबलता है)।

\textbf{20.} प्रति लीटर रस पर आयतन के हिसाब से लगभग 13\% ऐल्कोहॉल और
\qty{53}{L} गैस; और दोनों इसी से निकलते हैं कि खमीर, ऑक्सीजन के बिना, हर
ग्लूकोज़ का \omterm{def:g10:cell-metabolism:fermentation}{किण्वन} करके दो इथेनॉल और दो कार्बन डाइऑक्साइड बनाता है।
\end{solution}
